% 课程文档的统一字体配置。
% 正文与标题中的中文统一使用宋体，西文统一使用 Times New Roman。
% 代码环境保留等宽西文字体，以免破坏代码对齐。

\RequirePackage{fontspec}

% Overleaf 可在 assets/fonts/ 中放置合法授权的 Times New Roman 字体文件。
\IfFileExists{assets/fonts/times.ttf}{%
  \setmainfont{times.ttf}[
    Path=assets/fonts/,
    BoldFont=timesbd.ttf,
    ItalicFont=timesi.ttf,
    BoldItalicFont=timesbi.ttf
  ]
  \setsansfont{times.ttf}[
    Path=assets/fonts/,
    BoldFont=timesbd.ttf,
    ItalicFont=timesi.ttf,
    BoldItalicFont=timesbi.ttf
  ]
}{%
  \IfFontExistsTF{Times New Roman}{%
    \setmainfont{Times New Roman}
    \setsansfont{Times New Roman}
  }{%
    \setmainfont{TeX Gyre Termes}
    \setsansfont{TeX Gyre Termes}
    \PackageWarningNoLine{sduos-fonts}{Times New Roman unavailable; using TeX Gyre Termes}
  }%
}

% Overleaf 可在 assets/fonts/ 中放置合法授权的宋体字体文件。
\IfFileExists{assets/fonts/simsun.ttc}{%
  \setCJKmainfont{simsun.ttc}[
    Path=assets/fonts/,
    AutoFakeBold=2.5,
    AutoFakeSlant=0.2
  ]
  \setCJKsansfont{simsun.ttc}[
    Path=assets/fonts/,
    AutoFakeBold=2.5,
    AutoFakeSlant=0.2
  ]
  \setCJKmonofont{simsun.ttc}[
    Path=assets/fonts/,
    AutoFakeBold=2.5
  ]
}{%
  \IfFontExistsTF{SimSun}{%
    \setCJKmainfont{SimSun}[AutoFakeBold=2.5,AutoFakeSlant=0.2]
    \setCJKsansfont{SimSun}[AutoFakeBold=2.5,AutoFakeSlant=0.2]
    \setCJKmonofont{SimSun}[AutoFakeBold=2.5]
  }{%
    \setCJKmainfont{FandolSong-Regular}[
      BoldFont=FandolSong-Bold,
      AutoFakeSlant=0.2
    ]
    \setCJKsansfont{FandolSong-Regular}[
      BoldFont=FandolSong-Bold,
      AutoFakeSlant=0.2
    ]
    \setCJKmonofont{FandolSong-Regular}[BoldFont=FandolSong-Bold]
    \PackageWarningNoLine{sduos-fonts}{SimSun unavailable; using Fandol Song}
  }%
}
